%%%%%%%%%%%%%%%%%%%%%%%%%%%%%%%%%%%%%%%%%%%%%%%%%%%%%%%%%%%%%%%%%%%%%%%%%%%%%%%
%%
%%  Appendix.tex
%%  A Categorical Approach to Diffeology
%%  By Enxin Wu
%%
%%%%%%%%%%%%%%%%%%%%%%%%%%%%%%%%%%%%%%%%%%%%%%%%%%%%%%%%%%%%%%%%%%%%%%%%%%%%%%%

%%%%%%%%%%%%%%%%%%%%%%%%%%%%%%%%%%%%%%%%%%%%%%%%%%%%%%%%%%%%%%%%%%%%%%%%%%%%%%%
%%
%% MARK: Macros Enxin
%%
%%%%%%%%%%%%%%%%%%%%%%%%%%%%%%%%%%%%%%%%%%%%%%%%%%%%%%%%%%%%%%%%%%%%%%%%%%%%%%%

\newcommand{\uA}{\underline{A}}

\newtheorem{definition}{Definition}
\newtheorem{corollary}[definition]{Corollary}
\newtheorem{theorem}[definition]{Theorem}

\theoremstyle{remark}
\newtheorem{remark}[definition]{Remark}
\newtheorem{example}[definition]{Example}

\newcommand{\op}{\mathrm{op}}
\newcommand{\CCov}{\mathrm{Cov}}

\newcommand{\Set}{\mathfrak{Set}}
\newcommand{\sSet}{\mathfrak{sSet}}
\newcommand{\CCh}{\mathfrak{Ch}}

\renewcommand{\Top}{\mathfrak{Top}}
\renewcommand{\Diff}{\mathfrak{Diff}}

\newcommand{\Vect}{\mathfrak{Vect}}
\newcommand{\DVPB}{\mathfrak{Dvpb}}
\newcommand{\Pre}{\mathfrak{Pre}}
\newcommand{\Sh}{\mathfrak{Sh}}
\newcommand{\CCSh}{\mathfrak{CSh}}

\renewenvironment{proof}{\textit{proof}.}{~$\square$}

\def \infty{\inftyreg}

\newcommand{\colim}{\mathrm{colim}}

%%%%%%%%%%%%%%%%%%%%%%%%%%%%%%%%%%%%%%%%%%%%%%%%%%%%%%%%%%%%%%%%%%%%%%%%%%%%%%%
%%
%% MARK: Appendix
%%
%%%%%%%%%%%%%%%%%%%%%%%%%%%%%%%%%%%%%%%%%%%%%%%%%%%%%%%%%%%%%%%%%%%%%%%%%%%%%%%

\newpage
\markright{\ }
\
\newpage

%% Specific display for Enxin Wu Appendix
\MTversion{palatino}
\setcounter{footnote}{0}

\chapter*{Appendix: A Categorical Approach to Diffeology}
\vspace{-3ex}
\begin{center}
  Enxin Wu \\
  Department of Mathematics \\
  Shantou University, Guangdong \\
  P.R. China
\end{center}

\bigskip
From the very beginning,
diffeology was introduced from a geometric point of view.
But one can take an abstract categorical approach to understand and study this subject.
In this short appendix,
we briefly discuss this alternative aspect,
and see how to organize the basic theory of diffeology,
especially the categorical viewpoint of some geometric constructions and homotopy theory in diffeology.

A bit of category theory is needed to understand this appendix,
as the standard terminologies are not defined here.
The first section requires the basics of categories,
functors,
natural transformations,
limits,
colimits and cartesian closedness;
The next two sections use adjunctions and (left) Kan extensions.
Some familiarity with the corresponding contents in~\cite{McL71} suffices.

\section*{Diffeological spaces versus concrete sheaves}

The classical sheaf theory can be viewed as making the mathematical formalism of the local-to-global principle with underlying topological spaces.
For example, to define a suitable continuous map from a complicated topological space to $\RR$,
one could first decompose the space into simpler pieces (e.g., an open cover),
and define suitable continuous maps from these pieces to $\RR$.
Then one could try to patch them together to get a global one if it is possible.
This idea can be formalized rigorously using category theory as follows:

\begin{definition}
  For a topological space $X$,
  write $\cO_X$ for all open subsets of $X$,
  ordered by inclusions.
  A (set-valued) \textbf{presheaf} on $X$ is a functor $\cO_X^{\op} \to \Set$ from the opposite category of $\cO_X$ to the category $\Set$ of sets and functions.
  A (set-valued) \textbf{sheaf} on $X$ is a (set-valued) presheaf $\cF$ on $X$ such that for any $U \in \cO_X$ and any open covering $U = \bigcup_i U_i$,
  we have an equalizer diagram in $\Set$:
  \begin{center}
    \begin{tikzcd}[column sep=normal, row sep=normal, every label/.append style = {font = \small}]
      \cF(U) \arrow[r]  & \prod_i \cF(U_i) \arrow[r,shift left] \arrow[r, shift right] & \prod_{j,k} \cF(U_j \cap U_k)
    \end{tikzcd}
  \end{center}
  where the two parallel arrows are induced by the inclusions $U_j \cap U_k \to U_j$ and $U_j \cap U_k \to U_k$,
  respectively.
\end{definition}

It is easy to observe that if one abstracts the essential properties of $\cO_X$ out of the above definition,
one could define sheaves over other nice categories:

\begin{definition}
  A \textbf{presheaf} on an arbitrary category $\cA$ is a functor%
  \linebreak $\cA^{\op} \to \Set$.
  %  Write $\Pre(\cA)$ for the category of presheaves on $\cA$ and natural transformations between them.

  A \textbf{Grothendieck topology} $\CCov$ on a category $\cA$ is an assignment of a collection $\CCov(A)$ of collections of morphisms with the common codomain $A$ for each object $A$ of $\cA$ such that
  \begin{itemize}
    \item[*] $\{1_A:A \to A\} \in \CCov(A)$;
    \item[*] If $A' \to A$ is a morphism in $\cA$ and $\{A_i \to A\}_i \in \CCov(A)$,
    then each pullback $A' \times_A A_i$ exists in $\cA$,
    and $\{A' \times_A A_i \to A'\}_i \in \CCov(A')$;
    \item[*] If $\{A_i \to A\}_i \in \CCov(A)$,
    and $\{A_{ij} \to A_i\}_j \in \CCov(A_i)$ for each $i$,
    then $\{A_{ij} \to A_i \to A\}_{i,j} \in \CCov(A)$.
  \end{itemize}
  An element in $\CCov(A)$ is called a \textbf{covering} of $A$.

  A \textbf{site} is a category with a Grothendieck topology.

  A \textbf{sheaf} on a site $(\cA,\CCov)$ is a presheaf $\cF:\cA^{\op} \to \Set$ on $\cA$ such that for each covering $\{A_i \to A\}_{i \in I}$ of $A$,
  we have an equalizer diagram in $\Set$:
  \begin{center}
    \begin{tikzcd}[column sep=normal, row sep=normal, every label/.append style = {font = \small}]
      \cF(A) \arrow[r]  & \prod_{i \in I} \cF(A_i) \arrow[r,shift left] \arrow[r, shift right] & \prod_{j,k \in I} \cF(A_j \times_A A_k)
    \end{tikzcd}
  \end{center}
  where the two parallel arrows are induced by the canonical maps $A_j \times_A A_k \to A_j$ and $A_j \times_A A_k \to A_k$,
  respectively.
  %  We simply write $\Sh(\cA)$ for the category of sheaves over the site $(\cA,\CCov)$ and natural transformations between them.
\end{definition}

\begin{example}
  Let $X$ be a topological space,
  and let $\cO_X$ be the poset of all open subsets of $X$ together with inclusions.
  For any open subset $U$ of $X$,
  write $\CCov(U)$ for the collection of the usual open coverings of $U$.
  This gives rise to a Grothendieck topology on $\cO_X$,
  and the above two definitions of sheaves on $\cO_X$ coincide.
\end{example}

\begin{example}
  First example: Let $\cE^\infty$ be the category of all open subsets of $\RR^n$ for all $n \in \NN=\{0,1,2,3,\ldots\}$ and smooth (i.e., infinitely differentiable) maps between them.
  For any object $U$ of $\cE^\infty$,
  write $\CCov(U)$ for the collection of the usual open coverings of $U$.
  This gives rise to a Grothendieck topology on $\cE^\infty$.

  Second example: In a similar manner,
  if we only change the morphisms of $\cE^\infty$ to be $C^k$ (i.e., $k$-times continuously differentiable) for $k \in \NN$ or $C^\omega$ (i.e., analytic), then we again get sites denoted $\cE^k$ and $\cE^\omega$, respectively.

  Third example: Similarly, all open subsets of $\CC^n$ for all $n \in \NN$,
  all holomorphic maps between them,
  and the usual open coverings together form a site,
  denoted $\cC^\omega$.
\end{example}

From a differential geometer's point of view,
the general sheaves are abstract to work with.
It would be nice if there is an underlying set which dominates all the structures.
This leads to the following definition:

\begin{definition}
  A \textbf{concrete site} $(\cA,\CCov)$ is a site with a terminal object $1$ such that
  \begin{itemize}
    \item[*] the functor $\Hom_\cA(1,\cdot):\cA \to \Set$ is faithful;
    \item[*] for every object $A$ in $\cA$ and every covering $\{A_i \to A\}_i \in \CCov(A)$,
    the induced map $\coprod_i \Hom_\cA(1,A_i) \to \Hom_{\cA}(1,A)$ is surjective.
  \end{itemize}
  %
  A \textbf{concrete sheaf} $\cF$ over a concrete site $(\cA,\CCov)$ with a terminal object $1$ is a sheaf over $(\cA,\CCov)$ such that the natural map
  $$
    \cF(A) \to \Hom_{\Set}(\Hom_{\cA}(1,A),\cF(1))
  $$
  defined by $a \in \cF(A) \mapsto (f \in \Hom_{\cA}(1,A) \mapsto f^*(a) \in \cF(1))$,
  is injective for every object $A$ of $\cA$.

  We simply write $\CCSh(\cA)$ for the category of concrete sheaves over the concrete site $(\cA,\CCov)$ and natural transformations between them.
\end{definition}

\begin{example}
  All of the sites $\cE^\infty,\cE^k,\cE^\omega,\cC^\omega$ are concrete.
  They share a common terminal object which consists of a single point.
  The rest axioms for the concreteness are clear since every object (resp. morphism, covering) in each of these sites has a set-theoretical meaning.
\end{example}

\begin{example}
  The site $\cO_X$ is concrete if and only if every covering of any object $U$ in $\cO_X$ is essentially the trivial covering $1_U:U \to U$.
\end{example}

Here are some analogous definitions to that of a diffeological space and a smooth map between diffeological spaces:

\begin{definition}
  Given a concrete site $(\cA,\CCov)$ with a terminal object $1$,
  write $\uA$ for the set $\Hom_{\cA}(1,A)$ for each object $A$ of $\cA$.
  An \textbf{$\cA$-space} is a set $X$ together with a collection $\cS$ of maps (called the \textbf{structure maps}) $\uA \to X$ of sets for all objects $A$ of $\cA$ such that the following three axioms hold:
  \begin{itemize}
    \item[*] Every constant map $\uA \to X$ is in $\cS$;
    \item[*] If $\uA \to X \in \cS$ and $A' \to A$ is a morphism in $\cA$,
    then the composite $\uA' \to \uA \to X$ is in $\cS$;
    \item[*] Given a map $f:\uA \to X$,
    if there is a covering $\{A_i \to A\}_i$ of $A$ in $\CCov(A)$ such that each composite $\uA_i \to \uA \to X$ is in $\cS$,
    then so is the map $f$.
  \end{itemize}
  %
  A map $X \to Y$ between two $\cA$-spaces is called an \textbf{$\cA$-map},
  if it is a function of the two underlying sets,
  which sends every structure map of $X$ to a structure map of $Y$.
  $\cA$-spaces and $\cA$-maps form a category,
  denoted $\cA$-Spaces.
\end{definition}

\begin{theorem}
  For a concrete site $(\cA,\CCov)$,
  the two categories $\CCSh(\cA)$ and $\cA$-Spaces are equivalent.
\end{theorem}

\begin{proof}
  Let $1$ be a terminal object of the site.
  Define $F:\CCSh(\cA) \to \cA$-Spaces by sending a concrete sheaf $\cF$ over $\cA$ to $\cF(1)$;
  the structure maps on $\cF(1)$ consist of all $\cF(A)$ (viewed as a subset of $\Hom_{\Set}(\uA,\cF(1))$ from the definition of a concrete sheaf) for all objects $A$ in $\cA$;
  a morphism $\cF \to \cG$ is sent to the induced map $\cF(1) \to \cG(1)$.

  Define $G:\cA$-Spaces $\to \CCSh(\cA)$ by sending an $\cA$-space $X$ to a concrete sheaf $\cF_X$ such that $\cF_X(A) = $ the collection of structure maps of $X$ of the form $\uA \to X$,
  for every object $A$ of $\cA$; an $\cA$-map $f:X \to Y$ in $\cA$-Spaces is sent to a natural transformation $\cF_X \to \cF_Y$ induced by post-composition with $f$.

  By unraveling the definitions of a concrete sheaf and a natural transformations between two concrete sheaves,
  and the definitions of an $\cA$-space and an $\cA$-map between two $\cA$-spaces,
  one sees that $F$ and $G$ are functors which give rise to the equivalence of the two categories.%
  \footnote{We will need to use the fact that $\underline{A' \times_A A''} \cong \uA' \times_{\uA} \uA''$ in $\Set$ in the proof,
  and this fact can be easily verified by the universal property of pullback.}
\end{proof}

\begin{theorem}
  Let $(\cA,\CCov)$ be a concrete site.
  Then
  \begin{itemize}
    \item[*] $\cA$ is canonically a fully faithful subcategory of $\CCSh(\cA)$.
    \item[*] The category $\CCSh(\cA)$ is bicomplete and cartesian closed.
    \item[*] Every concrete sheaf over a concrete site $\cA$ is a colimit of a diagram of $\cA$ in $\CCSh(\cA)$.
  \end{itemize}
\end{theorem}

\begin{proof}
  We can replace $\CCSh(\cA)$ by the category $\cA$-Spaces.

  (1) Define $e:\cA \to \cA$-Spaces by sending an object $A$ of $\cA$ to an $\cA$-space $e(A)$ with the underlying set $\uA$. And $\uA' \to \uA$ is defined to be a structure map if and only if it is induced by a morphism $A' \to A$ in $\cA$.
  A morphism $A \to A''$ in $\cA$ then induces an $\cA$-map $\uA \to \uA''$.
  Therefore,
  $e$ is a functor.
  Moreover,
  two distinct morphisms $A \to A''$ in $\cA$ induces two distinct $\cA$-maps $e(A) \to e(A'')$ since $1_A:A \to A$ can distinguish them.
  Furthermore,
  every $\cA$-map $e(A) \to e(A'')$ is induced from a morphism $A \to A''$ in $\cA$,
  again by using the structure map of $e(A)$ induced by $1_A:A \to A$.
  These prove the first result.

  (2) Let $F:\cI \to \cA$-Spaces be a functor from a small category $\cI$.
  Write $U:\cA$-Spaces $\to \Set$ for the forgetful functor.
  Then $\lim F$ is an $\cA$-space which has $\lim (U \circ F)$ as the underlying set,
  and $\uA \to \lim F$ is a structure map if and only if each composite $\uA \to \lim F \to F(i)$ is a structure map of $F(i)$ for each object $i$ of $\cI$.
  Similarly, $\colim F$ is an $\cA$-space which has $\colim (U \circ F)$ as the underlying set.
  And $\uA \to \colim F$ is a structure map if and only if there is a covering $\{A_j \to A\}_j$ of $A$ in $\CCov(A)$ such that for each $j$ there is an object $i$ in $\cI$ and a structure map $\uA_j \to F(i)$ making the following diagram commutative:
  \begin{center}
    \begin{tikzcd}[column sep=normal, row sep=normal, every label/.append style = {font = \small}]
      {}              & {}                & F(i) \arrow[d] \\
      \uA_j \arrow[r] \arrow[urr] & \uA \arrow[r] & \colim F
    \end{tikzcd}
  \end{center}
  Given two $\cA$-spaces $X$ and $Y$,
  simply write $\cA(X,Y)$ for the collection of all morphisms $X \to Y$.
  Define $f:\uA \to \cA(X,Y)$ to be a structure map if and only if the adjoint $\tilde{f}:e(A) \times X \to Y$ is an $\cA$-map,
  where the domain is the product of two $\cA$-spaces and hence again an $\cA$-space as we have just proved that the category $\cA$-Spaces is complete.
  It is straightforward to check that $\cA(X,Y)$ with these structure maps is an $\cA$-space,
  and moreover,
  there is an adjoint pair of functors
  $$
    \cdot \times Y:\cA\text{-Spaces} \rightleftharpoons \cA\text{-Spaces}:\cA(Y,\cdot),
  $$
  for any fixed $\cA$-space $Y$.

  (3) Given an $\cA$-space $X$,
  write $\cA/X$ for the category whose objects are structure maps $\uA \to X$ and whose morphisms are commutative triangles
  \begin{center}
    \begin{tikzcd}[column sep=normal, row sep=normal, every label/.append style = {font = \small}]
      \uA' \arrow[rr,"\tilde f"] \arrow[dr, swap, "p"] & {} & \uA \arrow[dl,"q"] \\
      {}                & X &        {}
    \end{tikzcd}
  \end{center}
  where $p,q$ are structure maps of $X$ and $\tilde{f}$ is induced from a morphism $f:A' \to A$ in $\cA$.
  Let $\cA/X \to \cA$-Spaces be the functor sending the above commutative triangle to the $\cA$-map $e(A') \to e(A)$ induced by $f$.
  In other words, this functor is a composition $\cA/X \to \cA \to \cA$-Spaces. It is straightforward to check that the colimit of this functor is isomorphic to $X$ as $\cA$-spaces.
\end{proof}

\begin{corollary}
  First of all,
  $\CCSh(\cE^\infty)$ is equivalent to the category $\Diff$ of diffeological spaces.

  Then,
  $\Diff$ is bicomplete and cartesian closed.

  Next,
  every smooth manifold%
  \footnote{By a smooth manifold, we always require it to be finite-dimensional, Hausdorff, second-countable and without boundary.}
  is a gluing of Euclidean domains via diffeomorphisms,
  while every diffeological space is a gluing of Euclidean domains via smooth maps.
\end{corollary}
A bit more about concrete sheaves is discussed in~\cite{BH09}.

\section*{Constructions in diffeology revisited}

The general slogan is,
some structures of $\cE^\infty$ can be compatibly (hence functorially) pushed to diffeological spaces.
This slogan will be realized via Kan extension of a functor $\cE^\infty \to \cC$ (resp. $(\cE^\infty)^{\op} \to \cC$) to some category $\cC$ along the embedding $\cE^\infty \to \Diff$ (resp. $(\cE^\infty)^{\op} \to (\Diff)^{\op}$) as follows.

\textbf{D-topology.}
There is an adjoint pair of functors
$$
  D:\Diff \rightleftharpoons \Top:C
$$
where the functor $D$ sends a diffeological space $X$ to a topological space $D(X)$ with the same underlying set.
The topology (called the \textbf{D-topology}) on $D(X)$ consists of all subsets $A$ of $X$ such that $p^{-1}(A)$ is open in $U$ wi
